\documentclass[12pt,a4paper]{article}

\usepackage{fontspec}
\usepackage{polyglossia}
\setdefaultlanguage{russian}
\setotherlanguage{english}
\setmainfont{CMU Serif}
\setsansfont{CMU Sans Serif}
\setmonofont{CMU Typewriter Text}
\usepackage{geometry}
\geometry{top=2cm, bottom=2cm, left=2.5cm, right=2.5cm}
\usepackage{amsmath,amssymb,amsthm}
\usepackage{booktabs,longtable,float,xcolor}
\usepackage{hyperref}
\hypersetup{colorlinks, urlcolor=blue, citecolor=blue, linkcolor=black}
\usepackage{enumitem}
\sloppy
\tolerance=1000
\emergencystretch=3em
\numberwithin{equation}{section}

\title{\textbf{Модель MEML~75:}\\[0.3cm]
\large Феноменологическая параметризация фермионного сектора\\
Стандартной модели из геометрии $T^2/\mathbb{Z}_3$\\
с аксионоподобной (ALP) темной материей}
\author{Д.~В.~Юров}
\date{2026}

\begin{document}
\maketitle

\begin{abstract}
\noindent
Мы представляем модель MEML~75, которая параметризует фермионный сектор 
Стандартной модели через геометрию шестимерного орбифолда $T^2/\mathbb{Z}_3$ 
с тремя независимыми торами. Модель фитирует \textbf{13 фермионных наблюдаемых} 
(9 кварков $+$ 4 лептона) с помощью \textbf{16 свободных параметров} 
($N_{\rm dof} = 13-16 = -3$).

\textbf{Главное проверяемое предсказание:}
аксионоподобная темная материя с
$f_a = 1.9801\times10^{11}$~ГэВ,
$m_a = 28.407$~мкэВ, $\nu_a = 6.8688$~ГГц и ALP-фотонной связью
$g_{a\gamma\gamma} = +4.38\times10^{-15}$~ГэВ$^{-1}$ (DFSZ)
или $-1.13\times10^{-14}$~ГэВ$^{-1}$ (KSVZ), проверяемая
в ADMX/HAYSTAC (0.5--10~ГГц).

\textbf{Кварковый сектор.} Глобальный фит с функцией потерь
$\chi^2 = \sum_i ((O_i-O_i^{\rm target})/\sigma_i)^2$ дает
$\chi^2 = 6.195\times10^{-20}$ при
$\tau_u = -0.1577944085 + 2.3851782590\,i$,
$\tau_d = +0.3693905617 + 1.5410496057\,i$.

\textbf{Лептонный сектор.} Используя третий тор \(\tau_l\), комбинация $A_4$ модулярных форм $Y^{(2)}$ и $Y^{(4)}$ дает
$\chi^2 = 1.0065\times10^{-16}$ для наблюдаемых PMNS при
$\tau_l = 0.385257 + 0.958275\,i$. Модель предсказывает
$\theta_{12}=33.41^\circ$, $\theta_{13}=8.54^\circ$, $\theta_{23}=49.10^\circ$ и $J_{cp}=-0.018$.

\textbf{Вердикт:} модель \emph{феноменологическая}. Она интерполирует наблюдения фермионов с высокой точностью.
\end{abstract}

\newpage

\section{Введение}

Модулярные модели масс фермионов~\cite{Feruglio2017,Kobayashi2019}
строят иерархию масс и смешивание из геометрии компактификации.
В модели MEML~75 кварки и лептоны локализованы на трех
независимых торах $T^2/\mathbb{Z}_3$ с модулями $\tau_u$, $\tau_d$ и \(\tau_l\). Матрицы масс строятся из $A_4$ модулярных форм весов $(2,4,6)$.

\section{Геометрия $T^2/\mathbb{Z}_3$}
Орбифолд $T^2/\mathbb{Z}_3$ определяется как:
\begin{equation}
z = x + \tau y, \quad ds^2 = dx^2 + dx\,dy + dy^2,
\end{equation}
где $\tau = e^{i\pi/3} = \frac12 + i\frac{\sqrt3}{2}$. Действие $Z_3$ $z \to \omega z$ (где $\omega = e^{2\pi i/3}$) имеет три неподвижные точки, что дает топологическое объяснение трем поколениям.

\section{Объемная связь $\alpha_G$ (постулат)}
\begin{equation}
\alpha_G = \frac{2}{3\pi} = 0.2122065908.
\end{equation}
Геометрическая мотивация: $2/3$ --- доля объема фундаментальной области $\mathbb{Z}_3$-орбифолда; $1/\pi$ --- нормировка периода тэта-функции.

\section{Фактор подавления $\varepsilon$ (анзац)}
\begin{equation}
\varepsilon = \left(\frac{\alpha_G}{4\pi}\right)^2 = \frac{1}{36\pi^4} = 2.851662\times10^{-4}.
\end{equation}

\section{Условно выведенные величины}
\begin{align}
m_{KK} &= \frac{M_{Pl}^{\rm red}}{\sqrt{\mathcal{I}}\,e^{\pi^2}}
= 2.0838\times10^{14}~\text{ГэВ}, \\
M_6 &= \left(3 (M_{Pl}^{\rm red})^2 m_{KK}^2\right)^{1/4}
= 2.9646\times10^{16}~\text{ГэВ},
\end{align}
при $\mathcal{I} = 0.3653$.

\section{Аксионоподобный (ALP) сектор}
\subsection{Параметры ALP}
\begin{align}
f_a &= M_{Pl}^{\rm red}\,\varepsilon^2 = 1.9801\times10^{11}~\text{ГэВ}, \\
m_a &= \frac{(\chi^{1/4})^2}{f_a} = 28.407~\mu\text{эВ}, \\
\nu_a &= \frac{m_a c^2}{h} = 6.8688~\text{ГГц}.
\end{align}
Здесь $\chi^{1/4} \approx 75.0$~МэВ --- корень четвертой степени из топологической восприимчивости КХД. Мы подчеркиваем, что явная \(U(1)_{PQ}\) симметрия не вводится; объект является ALP.

\section{Кварковый сектор (11 параметров)}
\subsection{Матричная структура}
\begin{equation}
M_u = \begin{pmatrix}
a_u Y_1^{(2)} & b_u Y_3^{(4)} & c_u Y_2^{(6)} \\
a_u Y_3^{(2)} & b_u Y_2^{(4)} & c_u Y_1^{(6)} \\
a_u Y_2^{(2)} & b_u Y_1^{(4)} & c_u Y_3^{(6)}
\end{pmatrix}, \;
M_d = \begin{pmatrix}
a_1 Y_1^{(6)} & b Y_3^{(2)} & c Y_2^{(2)} \\
a_1 Y_3^{(6)} & b Y_2^{(2)} & c Y_1^{(2)} \\
a_2 Y_2^{(2)} & b Y_1^{(2)} & c Y_3^{(4)}
\end{pmatrix}.
\end{equation}

\subsection{Целевые значения}
\begin{table}[H]
\centering
\caption{Целевые значения (PDG~2024).}
\begin{tabular}{lc}
\toprule
Наблюдаемая & $O_i^{\rm target} \pm \sigma_i$ \\
\midrule
$m_u/m_t$ & $(1.27 \pm 0.30)\times10^{-3}/172.0$ \\
$m_c/m_t$ & $(0.620 \pm 0.020)/172.0$ \\
$m_d/m_b$ & $(2.70 \pm 0.20)\times10^{-3}/2.860$ \\
$m_s/m_b$ & $(55.0 \pm 5.0)\times10^{-3}/2.860$ \\
$m_b/m_t$ & $(2.860 \pm 0.030)/172.0$ \\
$|V_{us}|$ & $0.22501 \pm 0.00068$ \\
$|V_{cb}|$ & $0.04053 \pm 0.00042$ \\
$|V_{ub}|$ & $0.00382 \pm 0.00020$ \\
$|J|$ & $(3.080 \pm 0.130)\times10^{-5}$ \\
\bottomrule
\end{tabular}
\end{table}

\subsection{Результаты}
Независимая оптимизация дает $\chi^2 = 6.195\times10^{-20}$ при модулях, указанных в абстракте.

\section{Лептонный сектор (5 параметров)}
\label{sec:lepton}
Массы заряженных лептонов берутся диагональными (в ароматовом базисе), в то время как матрица масс нейтрино \(M_\nu\) строится из комбинации модулярных форм весов 2 и 4.

\subsection{Матричная структура}
\begin{equation}
M_\nu = g_1 M_{3_S}(Y^{(2)}) + g_2 M_{3_S}(Y^{(4)}) + d \cdot \mathbb{I},
\end{equation}
где триплетная свертка \(M_{3_S}(Y)\) имеет вид:
\begin{equation}
M_{3_S}(Y) = \begin{pmatrix}
2 Y_1 & -Y_3 & -Y_2 \\
-Y_3 & 2 Y_2 & -Y_1 \\
-Y_2 & -Y_1 & 2 Y_3
\end{pmatrix}.
\end{equation}
Диагональный член \(d \cdot \mathbb{I}\) нарушает \(\mu-\tau\) симметрию.

\subsection{Результаты}
Независимая оптимизация дает:
\begin{equation}
\boxed{\;\chi^2 = 1.0065\times10^{-16}
\quad\text{при}\quad
\tau_l = 0.385257 + 0.958275\,i,\;
g_1 = 9.9009,\; g_2 = 9.6201,\; d = -7.0956.\;}
\end{equation}

\begin{table}[H]
\centering
\caption{Наблюдаемые PMNS (NuFIT 5.2).}
\begin{tabular}{lccc}
\toprule
Наблюдаемая & Предсказание & Цель & Отклонение \\
\midrule
$\theta_{12}$ (солнечный) & $33.41^\circ$ & $33.41^\circ \pm 0.5^\circ$ & $\sim 10^{-9}$ \\
$\theta_{13}$ (реакторный) & $8.54^\circ$  & $8.54^\circ \pm 0.1^\circ$  & $\sim 10^{-9}$ \\
$\theta_{23}$ (атмосферный)& $49.10^\circ$ & $49.1^\circ \pm 0.5^\circ$  & $\sim 10^{-9}$ \\
$J_{cp}$ (CP-фаза)         & $-0.0180$     & $-0.018 \pm 0.005$          & $\sim 10^{-9}$ \\
\bottomrule
\end{tabular}
\end{table}

\section{Полный баланс}
\label{sec:balance}

\begin{table}[H]
\centering
\caption{Свободные параметры (16).}
\begin{tabular}{llc}
\toprule
Категория & Параметры & Количество \\
\midrule
Модули кварков & $\tau_u, \tau_d$ (комплексные) & 4 \\
Кварковые коэффициенты & $a_u, b_u, c_u, a_1, a_2, b, c$ & 7 \\
Модуль лептонов & \(\tau_l\) (комплексный) & 2 \\
Лептонные коэффициенты & $g_1, g_2, d$ & 3 \\
\midrule
\textbf{Всего} & & \textbf{16} \\
\bottomrule
\end{tabular}
\end{table}

\begin{equation}
N_{\rm dof} = 13 - 16 = -3 .
\end{equation}

\subsection{Информационные критерии}
\begin{table}[H]
\centering
\caption{AIC/BIC для фермионного сектора ($n=13$).}
\begin{tabular}{lccccc}
\toprule
Модель & $k$ & $\chi^2$ & AIC & BIC \\
\midrule
MEML~75 & 16 & $\approx 10^{-20}$ & $32.0$ & $36.7$ \\
SM       & 19 & $0.9$              & $38.9$ & $43.8$ \\
\bottomrule
\end{tabular}
\end{table}

\section{Заключение}
\textbf{Что было сделано:}
\begin{enumerate}
\item Совместный фит кварков и лептонов с 16 параметрами для 13 наблюдаемых ($N_{\rm dof} = -3$), плюс проверяемое предсказание для аксиона.
\item Кварковый сектор: $\chi^2 = 6.195\times10^{-20}$.
\item Лептонный сектор: $\chi^2 = 1.0065\times10^{-16}$ при \(\tau_l = 0.385257 + 0.958275\,i\). Нарушение \(\mu-\tau\) симметрии успешно достигается с помощью члена \(d \cdot \mathbb{I}\).
\item AIC/BIC с оговоркой $N_{\rm dof} < 0$.
\end{enumerate}
\textbf{Что не работает:} коэффициенты scale-вырождены; $\alpha_G$, $\varepsilon$, $\theta_i$ не выведены из первых принципов. Модель интерполирует наблюдения с помощью 16 параметров.

\appendix
\newpage

\section{$A_4$ модулярные формы: полные q-разложения}
Определение: $q = e^{2\pi i\tau}$.
\begin{align}
Y_1^{(2)} &= 1 + \sum_{n=1}^{17} C_1^{(n)} q^n, \\
Y_2^{(2)} &= -6\, q^{1/3}\left(1 + \sum_{n=1}^{11} C_2^{(n)} q^n\right), \\
Y_3^{(2)} &= -18\, q^{2/3}\left(1 + \sum_{n=1}^{10} C_3^{(n)} q^n\right),
\end{align}
$C_1 = [12, 36, 12, 84, 72, 36, 96, 180, 12, 216, 144, 36, 156, 96, 84, 240, 36]$;
$C_2 = [1, 7, 8, 18, 14, 30, 20, 42, 30, 56, 42]$;
$C_3 = [1, 2, 5, 6, 12, 8, 14, 10, 22, 16]$.

\textbf{Свертки:}
\begin{align}
Y_1^{(4)} &= (Y_1^{(2)})^2 - Y_2^{(2)} Y_3^{(2)}, \\
Y_2^{(4)} &= (Y_3^{(2)})^2 - Y_1^{(2)} Y_2^{(2)}, \\
Y_3^{(4)} &= (Y_2^{(2)})^2 - Y_1^{(2)} Y_3^{(2)},
\end{align}
\begin{align}
Y_1^{(6)} &= \tfrac13\left[2 Y_1^{(2)} Y_1^{(4)} - Y_2^{(2)} Y_3^{(4)} - Y_3^{(2)} Y_2^{(4)}\right], \\
Y_2^{(6)} &= \tfrac13\left[2 Y_3^{(2)} Y_3^{(4)} - Y_1^{(2)} Y_2^{(4)} - Y_2^{(2)} Y_1^{(4)}\right], \\
Y_3^{(6)} &= \tfrac13\left[2 Y_2^{(2)} Y_2^{(4)} - Y_3^{(2)} Y_1^{(4)} - Y_1^{(2)} Y_3^{(4)}\right].
\end{align}

\begin{thebibliography}{99}
\bibitem{Planck2018} Planck Collaboration, Astron. Astrophys. \textbf{641}, A6 (2020).
\bibitem{PDG2024} R.~L. Workman et al. (PDG), Prog. Theor. Exp. Phys. \textbf{2022}, 083C01 (2022).
\bibitem{ADMX2025} ADMX Collaboration, Phys. Rev. Lett. \textbf{134}, 021002 (2025).
\bibitem{Feruglio2017} F.~Feruglio, Eur. Phys. J. C \textbf{77}, 713 (2017).
\bibitem{Kobayashi2019} T.~Kobayashi et al., Prog. Theor. Exp. Phys. \textbf{2020}, 053B05 (2020).
\end{thebibliography}

\end{document}